\section{Hopf-Rinow 定理}\label{sec:hopf-rinow-theorem}

设 $(M,g)$ 是 Riemann 流形. 全正规邻域的存在性蕴含, 总可以用一条最短测地线连接两个充分接近的点 $p,q\in M$. 现在从整体上讨论同一问题.

\begin{RiemannExample}\label{ex:global-geodesics-of-euclidean-space-and-sphere}
\begin{enumerate}[label=(\arabic*),leftmargin=*]
\item 给定两个不同的点 $p,q\in\mathbb{R}^n$, 对 Euclidean 度量存在唯一一条连接这两点的测地线 (在相差重参数化的意义下).
\item 给定两个不同的点 $p,q\in S^n$, 对标准度量至少存在两条连接这两点的测地线, 它们互不为重参数化.
\item 若 $p\ne0$, 则对 $\mathbb{R}^n\setminus\{0\}$ 中的 Euclidean 度量, 不存在从 $p$ 到 $-p$ 的测地线.
\end{enumerate}
\end{RiemannExample}

在许多情形下 (例如 $\mathbb{R}^n\setminus\{0\}$), 存在不能延拓到参数的所有取值的测地线. 换言之, $\exp_p(v)$ 并非对所有 $v\in T_pM$ 都有定义.

\begin{RiemannDefinition}\label{def:geodesic-completeness}
若对每个点 $p\in M$, 映射 $\exp_p$ 都在 $T_pM$ 上有定义, 则称 Riemann 流形 $(M,\langle\cdot,\cdot\rangle)$ \textbf{测地完备}.
\end{RiemannDefinition}

连通 Riemann 流形还有另一种完备性概念, 源于任何这样的流形自然地是一个度量空间.

\begin{RiemannDefinition}\label{def:riemannian-distance}
设 $(M,\langle\cdot,\cdot\rangle)$ 是连通 Riemann 流形, 且 $p,q\in M$. $p$ 与 $q$ 之间的\textbf{距离}定义为
\[
  d(p,q)=\inf\{l(\gamma)\mid\text{$\gamma$ 是从 $p$ 到 $q$ 的分段可微曲线}\}.
\]
\end{RiemannDefinition}

注意, 若存在连接 $p$ 到 $q$ 的最短测地线 $c$, 则 $d(p,q)=l(c)$. 如下述命题所述, 函数 $d:M\times M\to[0,+\infty)$ 确实是一个距离.

\begin{RiemannProposition}\label{prop:riemannian-distance-defines-metric-topology}
$(M,d)$ 是度量空间, 即对所有 $p,q,r\in M$, $d$ 满足:
\begin{enumerate}[label=(\roman*),leftmargin=*]
\item \textbf{正性:} $d(p,q)\geq0$, 且 $d(p,q)=0$ 当且仅当 $p=q$;
\item \textbf{对称性:} $d(p,q)=d(q,p)$;
\item \textbf{三角不等式:} $d(p,r)\leq d(p,q)+d(q,r)$.
\end{enumerate}
$M$ 上诱导的度量空间拓扑与 $M$ 作为可微流形的拓扑一致.
\end{RiemannProposition}

\begin{RiemannProof}
见 Exercise~\ref{ex:hopf-rinow-theorem}(\ref{ex:proof-of-riemannian-metric-space-topology}).
\end{RiemannProof}

因此, 可以讨论作为度量空间的 $M$ 的完备性 (即 Cauchy 序列是否收敛). 完备与测地完备等价, 这就是下述定理的内容.

\begin{RiemannTheorem}[Hopf-Rinow]\label{thm:hopf-rinow}
设 $(M,\langle\cdot,\cdot\rangle)$ 是连通 Riemann 流形, $p\in M$. 以下断言等价:
\begin{enumerate}[label=(\roman*),ref=\roman*,leftmargin=*]
\item\label{thm:hopf-rinow-geodesically-complete} $M$ 测地完备;
\item\label{thm:hopf-rinow-metrically-complete} $(M,d)$ 是完备度量空间;
\item\label{thm:hopf-rinow-exponential-at-one-point} $\exp_p$ 在 $T_pM$ 上有定义.
\end{enumerate}
此外, 若 $(M,\langle\cdot,\cdot\rangle)$ 测地完备, 则对所有 $q\in M$, 存在连接 $p$ 到 $q$ 的测地线 $c$, 满足 $l(c)=d(p,q)$.
\end{RiemannTheorem}

\begin{RiemannProof}
显然 (\ref{thm:hopf-rinow-geodesically-complete}) $\Rightarrow$ (\ref{thm:hopf-rinow-exponential-at-one-point}).

首先证明若 (\ref{thm:hopf-rinow-exponential-at-one-point}) 成立, 则对所有 $q\in M$, 存在连接 $p$ 到 $q$ 的测地线 $c$, 满足 $l(c)=d(p,q)$. 设 $d(p,q)=\rho$. 若 $\rho=0$, 则 $q=p$, 无需证明. 若 $\rho>0$, 取 $\varepsilon\in(0,\rho)$ 使 $S_\varepsilon(p)$ 是一个正规球面 (它是 $M$ 的一个紧子流形). 于是连续函数 $x\mapsto d(x,q)$ 有一个最小值点 $x_0\in S_\varepsilon(p)$. 此外, $x_0=\exp_p(\varepsilon v)$, 其中 $\lVert v\rVert=1$. 考虑测地线 $c_v(t)=\exp_p(tv)$. 将证明 $q=c_v(\rho)$. 为此, 考虑集合
\[
  A=\{t\in[0,\rho]\mid d(c_v(t),q)=\rho-t\}.
\]
由于映射 $t\mapsto d(c_v(t),q)$ 连续, $A$ 是闭集. 此外, $A\ne\varnothing$, 因为显然 $0\in A$. 现在证明任何点 $t_0\in[0,\rho)$ 都不可能是 $A$ 的最大值. 这蕴含 $A$ 的最大值必为 $\rho$ (因而 $d(c_v(\rho),q)=0$, 即 $c_v(\rho)=q$, 所以 $c_v$ 连接 $p$ 到 $q$, 且 $l(c_v)=\rho$). 取 $t_0\in A\cap[0,\rho)$, $r=c_v(t_0)$, 并取 $\delta\in(0,\rho-t_0)$ 使 $S_\delta(r)$ 为正规球面. 设 $y_0$ 是连续函数 $y\mapsto d(y,q)$ 在紧集 $S_\delta(r)$ 上的一个最小值点 (Fig.~\ref{fig:hopf-rinow-proof}). 将证明 $y_0=c_v(t_0+\delta)$. 事实上, 有
\RGfigure[.55\textwidth]{hopf-rinow-proof}{Hopf-Rinow 定理的证明}
\[
  \rho-t_0=d(r,q)=\delta+\min_{y\in S_\delta(r)}d(y,q)=\delta+d(y_0,q),
\]
所以
\begin{equation}\label{eq:distance-from-minimum-on-normal-sphere}
  d(y_0,q)=\rho-t_0-\delta.
\end{equation}
由三角不等式得到
\[
  d(p,y_0)\geq d(p,q)-d(y_0,q)=\rho-(\rho-t_0-\delta)=t_0+\delta,
\]
而沿 $c_v$ 从 $p$ 到 $r$, 再沿测地线从 $r$ 到 $y_0$ 的分段可微曲线, 长度为 $t_0+\delta$. 因此它是最短曲线, 从而是一条 (重参数化的) 测地线. 如前所述, 得 $y_0=c_v(t_0+\delta)$. 因此, 由 \eqref{eq:distance-from-minimum-on-normal-sphere} 可写成
\[
  d(c_v(t_0+\delta),q)=\rho-(t_0+\delta),
\]
这意味着 $t_0+\delta\in A$, 所以 $t_0$ 不可能是 $A$ 的最大值.

现在证明 (\ref{thm:hopf-rinow-exponential-at-one-point}) $\Rightarrow$ (\ref{thm:hopf-rinow-metrically-complete}). 先证明任何有界闭子集 $K\subset M$ 都紧. 事实上, 若 $K$ 有界, 则对某个 $R>0$, 有 $K\subset B_R(p)$, 其中
\[
  B_R(p)=\{q\in M\mid d(p,q)<R\}.
\]
如前所见, $p$ 可由长度小于 $R$ 的测地线连接到 $B_R(p)$ 中任意一点, 因而 $B_R(p)\subset\exp_p\bigl(\overline{B_R(0)}\bigr)$. 由于 $\exp_p:T_pM\to M$ 连续且 $\overline{B_R(0)}$ 紧, 集合 $\exp_p\bigl(\overline{B_R(0)}\bigr)$ 也紧. 因此 $K$ 是紧集的闭子集, 从而 $K$ 紧. 现在, 若 $\{p_n\}$ 是 $M$ 中的一个 Cauchy 序列, 则其闭包有界, 从而紧. 因此 $\{p_n\}$ 必有收敛子序列, 所以它本身必收敛.

最后证明 (\ref{thm:hopf-rinow-metrically-complete}) $\Rightarrow$ (\ref{thm:hopf-rinow-geodesically-complete}). 设 $c$ 是一条对 $t<t_0$ 有定义的测地线. 不失一般性, 可假设它已\textbf{归一化}, 即 $\lVert\dot c(t)\rVert=1$. 设 $\{t_n\}$ 是收敛到 $t_0$ 的递增实数序列. 由于 $d(c(t_m),c(t_n))\leq|t_m-t_n|$, 可知 $\{c(t_n)\}$ 是一个 Cauchy 序列. 因假设 $M$ 完备, 可知 $c(t_n)\to p\in M$, 且容易看出, 当 $t\to t_0$ 时, $c(t)\to p$. 设 $B_\varepsilon(p)$ 为以 $p$ 为中心的正规球. 则 $c$ 在此正规球内可以延拓到 $t_0$ 之后.
\end{RiemannProof}

\begin{RiemannCorollary}\label{cor:compact-riemannian-manifold-geodesically-complete}
若 $M$ 紧, 则 $M$ 测地完备.
\end{RiemannCorollary}

\begin{RiemannProof}
任何紧度量空间都完备.
\end{RiemannProof}

\begin{RiemannCorollary}\label{cor:closed-submanifold-complete-in-induced-metric}
若 $M$ 是完备连通 Riemann 流形的闭连通子流形, 并赋予诱导度量, 则 $M$ 完备.
\end{RiemannCorollary}

\begin{RiemannProof}
设 $M$ 是完备连通 Riemann 流形 $N$ 的闭连通子流形. 设 $d$ 是 $N$ 上度量确定的距离, $d^*$ 是 $M$ 上诱导度量确定的距离. 因为 $M$ 上的任何曲线也是 $N$ 上的曲线, 所以 $d\leq d^*$. 设 $\{p_n\}$ 是 $(M,d^*)$ 中的一个 Cauchy 序列. 则 $\{p_n\}$ 是 $(N,d)$ 中的 Cauchy 序列, 因而在 $N$ 中收敛到一点 $p\in M$ (因为 $N$ 完备且 $M$ 闭). 由于 $M$ 的拓扑由 $N$ 的拓扑诱导, 可知在 $M$ 上有 $p_n\to p$.
\end{RiemannProof}

\begin{RiemannExercise}\label{ex:hopf-rinow-theorem}
\begin{enumerate}[label=(\arabic*),ref=\arabic*,leftmargin=*]
\item\label{ex:proof-of-riemannian-metric-space-topology} 证明 Proposition~\ref{prop:riemannian-distance-defines-metric-topology}.
\item\label{ex:metric-ball-in-slit-plane} 考虑带 Euclidean 度量的 $\mathbb{R}^2\setminus\{(x,0)\mid-3\leq x\leq3\}$. 求 $B_7(0,4)$.
\item\label{ex:heine-borel-property-and-completeness}
\begin{enumerate}[label=(\alph*),leftmargin=*]
\item 证明连通 Riemann 流形完备当且仅当紧集就是有界闭集.
\item 给出一个含有非紧的有界闭集的连通 Riemann 流形的例子.
\end{enumerate}
\item\label{ex:completeness-of-homogeneous-manifolds} 若对任意两点 $p,q\in M$, 都存在等距映射 $f:M\to M$ 使 $f(p)=q$, 则称 Riemann 流形 $(M,\langle\cdot,\cdot\rangle)$ \textbf{齐性}. 证明:
\begin{enumerate}[label=(\alph*),leftmargin=*]
\item 任意齐性 Riemann 流形都完备;
\item 若 Lie 群 $G$ 容许双不变度量 [参见 Exercise~\ref{ex:minimizing-properties-of-geodesics}(\ref{ex:bi-invariant-metrics-on-lie-groups})], 则指数映射 $\exp:\mathfrak{g}\to G$ 是满射;
\item $SL(2,\mathbb{R})$ 不容许双不变度量.
\end{enumerate}
\item\label{ex:ambrose-and-cartan-hadamard-theorems} 设 $(M,g)$ 是完备 Riemann 流形. 证明:
\begin{enumerate}[label=(\alph*),leftmargin=*]
\item (\emph{Ambrose 定理}) 若 $(N,h)$ 是 Riemann 流形, $f:M\to N$ 是局部等距映射, 则 $f$ 是覆盖映射;
\item 存在不是覆盖映射的满局部等距映射;
\item (\emph{Cartan–Hadamard 定理}) 若 $(M,g)$ 具有非正曲率 [参见 Exercise~\ref{ex:minimizing-properties-of-geodesics}(\ref{ex:jacobi-fields-and-cut-locus})], 则对每个点 $p\in M$, 指数映射 $\exp_p:T_pM\to M$ 是覆盖映射.
{\small (\textbf{Remark:} 特别地, 若 $M$ 单连通, 则 $M$ 必与 $\mathbb{R}^n$ 微分同胚).}
\end{enumerate}
\end{enumerate}
\end{RiemannExercise}
